% REARIT-ID: REARIT-PXXX
% REARIT-VERSION: 0.1.0
% REARIT-STATUS: PROPOSTO
% REARIT-TITLE: Título do princípio
% REARIT-DATE: YYYY-MM-DD
% REARIT-ORIGIN: contexto / evento reflexivo
% REARIT-SUPERSEDES:
% REARIT-SUPERSEDED-BY:
%
% Relações genealógicas, quando presentes, usam chave histórica completa:
% REARIT-PXXX@MAJOR.MINOR.PATCH

\chapter{REARIT-PXXX --- Título do princípio}
\section{Natureza desta formulação}
Preservar a genealogia sem projetar retrospectivamente a formulação atual.
\section{Enunciado normativo}
Enunciar o princípio de forma verificável.
\section{Relação com REA}
Explicar perceber, interpretar e revisar.
\section{Relação com RIT}
Explicar identidade retrospectiva, plasticidade prospectiva, proveniência e
legitimidade da transição.
\section{Invariantes decorrentes}
Enumerar invariantes testáveis.
\section{Consequências de engenharia}
Explicitar consequências sem confundi-las com o princípio.
\section{Aplicações conhecidas}
Registrar aplicações factuais.
\section{Critérios de teste}
Definir como uma aplicação pode demonstrar aderência.
\section{Hipóteses abertas}
Preservar perguntas ainda não resolvidas.
\section{Histórico de evolução}
Registrar revisões, deprecações e sucessores.
